\chapter{The Platonic Variation and RH Completion}
\label{v4-arith:platonic-rh-completion}

\emph{The arithmetic branch's correct geometric variation is an
operadic adelic surface: holomorphic at \(\infty\), Frobenius-oriented
at finite primes, and glued by the pointed \(E_{2}\) restricted tensor
product of Iwahori Hecke categories: not a scheme-theoretic surface,
not a literal almost-complex manifold at finite primes. The upgraded
object \(\mathfrak{S}^{\mathrm{Ar}}_{2}\) drives the Riemann Hypothesis
completion inside Vol~IV as a conditional theorem.}

\section{The Right Variation}
\label{v4-arith:platonic:right-variation}

\begin{definition}[Platonic \(2\)-arithmetic surface]
\label{v4-arith:platonic:def:S2}
A \emph{Platonic \(2\)-arithmetic surface}
for the primary arithmetic branch is the stratified adelic-diamond
factorisation datum
\begin{equation}
\label{v4-arith:platonic:eq:S2}
\mathfrak{S}^{\mathrm{Ar}}_{2}
\;:=\;
\bigl(|\overline{C}|,\{S_{v},\mathfrak{J}^{\mathrm{Ar}}_{v},\mathcal{C}_{v},\eta_{v},u_{v}\}_{v},
\bigotimes_{v\in|\overline{C}|}^{\prime}(\mathcal{C}_{v},u_{v})\bigr),
\end{equation}
where:
\begin{itemize}
\item \(|\overline{C}|=\mathbb{P}\sqcup\{\infty\}\) is the set of
places of \(\overline{\mathrm{Spec}\,\mathbb{Z}}\);
\item \(S_{\infty}=\mathbb{C}_{\infty}\),
\(\mathfrak{J}^{\mathrm{Ar}}_{\infty}=i\), and
\(\mathcal{C}_{\infty}\) is the archimedean holomorphic
factorisation category with Harish--Chandra
\(K_{\infty}\)-spherical unit \(u_{\infty}\);
\item for \(p<\infty\), \(S_{p}\) is the Fargues--Scholze
perfectoid-diamond Hecke chart, \(\mathfrak{J}^{\mathrm{Ar}}_{p}\)
is the Frobenius-semilinear untilt-orientation with Adams
operation \(\psi_{p}\) and prismatic/Nygaard comparison datum, and
\(\mathcal{C}_{p}=\mathrm{Hk}^{\mathrm{Iw}}_{GL_{2},p}\) is the full
non-semisimple Iwahori Hecke category with spherical unit
\(u_{p}=|0\rangle_{p}\);
\item \(\eta_{v}\colon\mathrm{Sp}\to\mathcal{C}_{v}\) is the pointed
\(E_{2}\)-algebra map with image \(u_{v}\); at almost every finite place
\(u_{v}\) is the monoidal unit, and at the remaining places it is a
compact central idempotent algebra object;
\item the tensor product in~\eqref{v4-arith:platonic:eq:S2} is the
Tate restricted tensor product of pointed \(E_{2}\)-monoidal stable
presentable \(\infty\)-categories in the sense of
Section~\ref{v4-arith:tate:setup:data}.
\end{itemize}
The word ``surface'' refers to the operadic dimension: the global
category is \(E_{2}\)-monoidal, not a smooth scheme-theoretic surface
over \(\mathbb{Z}\).
\end{definition}

\begin{remark}[almost complex versus arithmetic complex]
\label{v4-arith:platonic:rem:almost-complex}
The symbol \(\mathfrak{J}^{\mathrm{Ar}}\) is literal almost-complex
geometry only on \(S_{\infty}=\mathbb{C}_{\infty}\). At a finite prime
it is the Frobenius-oriented substitute: a semilinear orientation of
the diamond chart together with Adams and prismatic comparison data.
The finite place supplies no tangent-bundle endomorphism satisfying
\(J^{2}=-1\); its role is to make the \(E_{2}\) Iwahori Hecke datum
functorial.
This convention governs every occurrence of
\(\mathfrak{J}^{\mathrm{Ar}}\) in the arithmetic branch.
\end{remark}

\begin{remark}[source anchors and dependency separation]
\label{v4-arith:platonic:rem:sources}
The finite charts use Scholze diamonds and Fargues--Scholze local
Langlands geometry as listed in
Section~\ref{v4-arith:app:A:langlands} and
Section~\ref{v4-arith:app:A:prismatic}. The
archimedean factorisation \(D\)-module fibre is the
Beilinson--Drinfeld/Francis--Gaitsgory/Costello--Gwilliam fibre of
Section~\ref{v4-arith:app:A:beilinson-etc} and
Section~\ref{v4-arith:app:A:costello}. The
arithmetic \(\mathcal{D}^{\dagger}\) and overconvergent
\(F\)-isocrystal language is only a local coefficient shadow, with
primary sources recorded in
Section~\ref{v4-arith:app:A:arith-Ddagger}. The definition of
\(\mathfrak{S}^{\mathrm{Ar}}_{2}\) is complete on the data above; F3
and the wild-prime BZN comparison enter only when identifying
\(\mathcal{V}^{\mathrm{prim}}\).
\end{remark}

\begin{question}[Comparison of arithmetic base data]
\label{v4-arith:platonic:thm:base-upgrade}
Construct a comparison from the base and coefficient system of
Definition~\ref{v4-arith:def:aCFT} to the specified datum
\(\mathfrak S^{\mathrm{Ar}}_2\), preserving the archimedean
factorization, the finite-place collision maps and their
Adams/Frobenius operations, and the global tensor structure.
\end{question}

The coefficient identity of
Theorem~\ref{v4-arith:thm:arch-vs-finite-asymmetry} does not construct
the finite-place maps. A pointing gives an \(E_0\)-datum;
the failure of a prescribed unary vacuum to extend in
Example~\ref{v4-ex:word-vacuum-obstruction} shows why its geometric
extension requires further operations. Once the local
\(E_2\)-monoidal categories and their compatible unital transitions
are supplied, Corollary~\ref{v4-arith:tate:cor:F2-closed} transports
that structure to their restricted tensor product. Its hypotheses
remain part of the proposed arithmetic comparison.

\section{Factorisable Arithmetic Modules}
\label{v4-arith:platonic:fact-modules}

The local categories \(\mathcal C_v\) of
\(\mathfrak S^{\mathrm{Ar}}_2\), when supplied, have a finite
tensor assembly over the place labels. Their position diagrams and
geometric comparison maps remain part of the local data.

\begin{definition}[factorisable arithmetic crystal]
\label{v4-arith:platonic:def:fact-crystal}
A \emph{factorisable arithmetic
crystal} on \(\mathfrak{S}^{\mathrm{Ar}}_{2}\) is a functor
\begin{equation}
\label{v4-arith:platonic:eq:fact-crystal}
\mathscr{D}^{\mathrm{Ar}}\colon
\mathrm{Fin}^{u}(|\overline{C}|)
\longrightarrow
\mathrm{Alg}_{\mathcal{O}}(\mathrm{Pr}^{\mathrm{L}}_{\mathrm{st}})
\end{equation}
such that
\[
\mathscr{D}^{\mathrm{Ar}}(S)\simeq\bigotimes_{v\in S}\mathcal{C}_{v},
\qquad
\mathscr{D}^{\mathrm{Ar}}(S_{1}\sqcup S_{2})
\simeq
\mathscr{D}^{\mathrm{Ar}}(S_{1})\otimes
\mathscr{D}^{\mathrm{Ar}}(S_{2})
\]
for disjoint finite subsets \(S_{1},S_{2}\subset|\overline{C}|\),
with coherent transition maps given by insertion of the specified
spherical units \(u_v\). The category \(\mathrm{Fin}^{u}\) is
defined in Definition~\ref{v4-arith:tate:def:factorization-discrete}; its
arrows add only places with supplied unit maps. An arithmetic
realization also specifies a holomorphic factorization model at
infinity and local geometric models at finite places, with maps
to their coefficient categories. Proposed finite models include
Iwahori Hecke categories on diamond charts and spectral
ind-coherent categories with nilpotent singular support.
An overconvergent \(F\)-isocrystal or arithmetic
\(\mathcal D^{\dagger}\) presentation on a chosen smooth formal
chart requires its own comparison. The global identification with
\(\mathrm{IndCoh}_{\mathrm{Nilp}}(\mathcal{X}^{\mathrm{Sel}}_{\bar\rho,\det,\mathcal{L}})\)
remains the F1/LACD/BC-diamond-glue comparison, not a definitional
equivalence.
\end{definition}

\begin{proposition}[Finite tensor assembly]
\label{v4-arith:platonic:prop:Dmodule-upgrade}
Given the \(\mathcal O\)-monoidal coefficient categories and
compatible unit maps of Setup~\ref{v4-arith:tate:setup:data},
the assignment \(S\mapsto\bigotimes_{v\in S}\mathcal C_v\)
has the discrete factorization and transition maps required in
Definition~\ref{v4-arith:platonic:def:fact-crystal}.
\end{proposition}

\begin{proof}
Apply Theorem~\ref{v4-arith:tate:thm:BD-automatic-discrete}.
Its tensor associators give the disjoint-set equivalences, and
its covariant unit insertions give the required transitions.
\end{proof}

This construction does not identify a geometric factorization
\(D\)-module with a finite-place coefficient category. The
holomorphic, diamond, spectral and differential presentations
require their stated comparison maps. At \(p=2,3\) no
semisimple truncation is intended. To include a wild contribution
in \(\mathcal C_p\), the actual local category and trace comparison
of Question~\ref{v4-arith:bzn-rcompletion:thm:F3-BZN-Ar-corrected}
must be supplied. Over a common field, assembling the proposed
local traces also includes the mixed degree-two term of
Lemma~\ref{v4-arith:wild:four-term-tensor}; an additive degree-one
formula alone is not the full tensor.

\section{The Arithmetic-Chiral Bar Complex on the Variation}
\label{v4-arith:platonic:bar}

The ordinary arithmetic complexes of
Theorem~\ref{v4-arith:pos-costello:thm:C2-strict-model} have explicit
square-zero differentials. The arithmetic-chiral comparison additionally
requires common coefficients, local chiral operations and residue couplings.

\begin{definition}[Finite-stage bar data]
\label{v4-arith:platonic:def:bar}
Fix a commutative ring \(R\) and a filtered poset \(\mathcal I\) of finite
sets of places containing \(\infty\). For each \(S\in\mathcal I\), give
cochain complexes \(A_\infty,G_S,A_p\), \(p\in S_f\), over \(R\).
Their intended chiral, residue and local arithmetic realizations are
additional data. Use algebraic tensor products and direct-sum totalization:
\begin{equation}
\label{v4-arith:platonic:eq:finite-bar}
M_S=A_\infty\otimes_R G_S\otimes_R\bigotimes_{p\in S_f}A_p.
\end{equation}
Order the factors as \(\infty,\mathrm{PG},p_1,\ldots,p_r\).
On a homogeneous tensor, let \(\widetilde d_a\) apply the differential
in factor \(a\) with sign \((-1)^{\sum_{b<a}|x_b|}\). Set
\begin{equation}
\label{v4-arith:platonic:eq:bar-differential}
D_{0,S}=\sum_a\widetilde d_a,\qquad D_S=D_{0,S}+\tau_S,
\end{equation}
where the specified degree-one operator
\(\tau_S\in\operatorname{End}_R^1(M_S)\) satisfies
\begin{equation}
\label{v4-arith:platonic:eq:bar-local-differentials}
D_{0,S}\tau_S+\tau_SD_{0,S}+\tau_S^2=0.
\end{equation}
Require degree-zero transition maps \(j_{ST}\) with
\(j_{SS}=1\), \(j_{TU}j_{ST}=j_{SU}\), and
\[
j_{ST}D_{0,S}=D_{0,T}j_{ST},\qquad j_{ST}\tau_S=\tau_Tj_{ST}.
\]
The resulting algebraic finite-support complex is
\(B^{\mathrm{Ar}}=\operatorname*{colim}_{S\in\mathcal I}(M_S,D_S)\).
Its arithmetic-chiral interpretation requires identification of these
carriers and operators with the intended arithmetic and chiral objects.
\end{definition}

\begin{lemma}[Insertion of closed units]
\label{v4-arith:platonic:lem:closed-units}
If \(G_S\to G_T\) is a cochain map and each new factor has
\(u_p\in A_p^0\) with \(d_pu_p=0\), insertion commutes with \(D_0\).
It commutes with \(D\) when it also intertwines \(\tau\).
\end{lemma}
\begin{proof}
An inserted factor contributes \(d_pu_p=0\), and its degree zero
leaves the other tensor signs unchanged. The remaining difference is
\(D_Tj_{ST}-j_{ST}D_S=\tau_Tj_{ST}-j_{ST}\tau_S\).
\end{proof}

\begin{proposition}[Finite-stage differential]
\label{v4-arith:platonic:prop:finite-stage-strictification}
The data of Definition~\ref{v4-arith:platonic:def:bar} give
\(D_S^2=0\) and cochain transition maps over any commutative ring \(R\),
including rings in which \(2\) is not invertible.
\end{proposition}
\begin{proof}
Each \(\widetilde d_a^2=0\). For \(a<b\), the first differential
increases by one the exponent in the sign of the second, so
\(\widetilde d_a\widetilde d_b+\widetilde d_b\widetilde d_a=0\).
Thus \(D_{0,S}^2=0\), and expansion gives
\[
D_S^2=D_{0,S}^2+D_{0,S}\tau_S+\tau_SD_{0,S}+\tau_S^2=0.
\]
The two transition identities imply \(D_Tj_{ST}=j_{ST}D_S\).
\end{proof}

This operator identity involves every summand of \(D_0\).
Scalar reciprocity for residue homomorphisms does not establish it.
No coalgebra structure, convolution product or degree-one tame twisting
operator is implicit. The ordinary complexes \(C_{r,s}\) of
Theorem~\ref{v4-arith:pos-costello:thm:C2-strict-model} give a concrete
untwisted example: \(D_0=I+J+L\), \(\tau=0\). Their mixed-square
identities, norm maps, coefficient reductions and completed exterior
comparison are proved there; their identification with a tame-weighted
chiral bar remains a separate construction.

\begin{proposition}[Comparison of finite-stage systems]
\label{v4-arith:platonic:prop:bar-comparison}
Compatible cochain maps \(q_S:N_S\to M_S\) between filtered systems
of cochain complexes over \(R\) induce a cochain map on their colimits.
If every \(q_S\) is a quasi-isomorphism, so is the induced map.
\end{proposition}
\begin{proof}
Transition compatibility defines the map. Exactness of filtered colimits
of \(R\)-modules gives
\(H^n(\operatorname{colim}q_S)=\operatorname{colim}H^n(q_S)\).
\end{proof}

Application to Definition~\ref{v4-arith:def:arith-bar-concrete} requires
actual common coefficient realizations and stage maps, differential and
transition identities, and quasi-isomorphism proofs. Finite support
alone does not identify a single Iwasawa action with the tensor of
local Iwasawa actions.

\begin{theorem}[Closure for specified arithmetic-chiral data]
\label{v4-arith:platonic:thm:bar-closure}
Suppose the carriers, operators and transitions of
Definition~\ref{v4-arith:platonic:def:bar} are constructed on the intended
arithmetic-chiral domain. Their algebraic finite-support colimit is a
cochain complex. For a completed version, suppose complete Hausdorff
filtrations are specified and the operators and transitions descend
compatibly to every quotient. Their inverse-limit extensions also have
square zero and commute with transitions.
\end{theorem}
\begin{proof}
Proposition~\ref{v4-arith:platonic:prop:finite-stage-strictification}
proves the finite-stage identities and hence the colimit assertion.
The completed identities hold in each quotient and therefore in the
separated inverse limit.
\end{proof}

The tame-weighted arithmetic realization is the construction problem of
Question~\ref{v4-arith:pos-costello:thm:C2-closed}. The common arithmetic
and real coefficient complex and its comparison maps are requested in
Question~\ref{v4-arith:pos-costello:thm:C4-closed}; the analytic carrier
and state comparison are requested in
Question~\ref{v4-arith:pos-costello:thm:C3-closed}. The ordinary bar and
integral oscillator models supply local constructions for these
questions. An algebraic colimit is not thereby identified with a
completed restricted tensor product.

\section{RH Completion}
\label{v4-arith:platonic:RH-completion}

\begin{definition}[RH completion hypothesis]
\label{v4-arith:platonic:def:RH-package}
The \emph{RH completion hypothesis}
is the conjunction of:
\begin{itemize}
\item the Platonic surface \(\mathfrak{S}^{\mathrm{Ar}}_{2}\) of
Definition~\ref{v4-arith:platonic:def:S2};
\item an arithmetic-chiral realization of Definition~\ref{v4-arith:platonic:def:bar},
including common coefficients, degree-one operators, transition
identities and the required completed comparison maps;
\item the arithmetic Heisenberg Koszul biconditional, conditional on
the arithmetic Positselski completion input of
Corollary~\ref{v4-arith:pos-costello:cor:arith-Heis-Koszul};
\item the arithmetic Petersson positive-form condition, conditional on
the restricted Koszul--Petersson compatibility input recorded in
Corollary~\ref{v4-arith:kp:cor:P3-unconditional};
\item H-P\(^{\mathrm{Ar}}\) in the equivalent determinant-trace
form of Theorem~\ref{v4-arith:hpvbconv:thm:HPAr-equivalence};
\item VB\(^{*}\) on the specified arithmetic carrier. It may be
supplied by (Reg), meaning positivity, determinacy and the required
spectral support of the identified arithmetic moment functional,
together with item~1 of
Hypothesis~\ref{v4-arith:hpvbconv:hyp:operator-quantitative-comparison}.
Substitution of (HB+) or (BD+) additionally requires the corresponding
operator or quantitative comparison from that hypothesis.
\end{itemize}
\end{definition}

\begin{theorem}[Platonic RH completion, conditional]
\label{v4-arith:platonic:thm:RH-completion}
Conditional on the RH completion
hypothesis of Definition~\ref{v4-arith:platonic:def:RH-package}, the Riemann
Hypothesis holds.
\end{theorem}

\begin{proof}
The surface \(\mathfrak{S}^{\mathrm{Ar}}_{2}\) and the bar complex
\(B^{\mathrm{Ar}}(\mathcal{V}^{\mathrm{prim}})\) of
Definition~\ref{v4-arith:platonic:def:bar} place
\(\mathcal{V}^{\mathrm{prim}}\) (a Hochschild trace class; see
Remark~\ref{v4-arith:rh:rem:equivalence-to-RH}) in the correct
arithmetic-chiral bar framework. The conditional Koszul input gives
P1 in the abelian Heisenberg domain
(Corollary~\ref{v4-arith:pos-costello:cor:arith-Heis-Koszul}); P2 is Tate's
functional-equation/Zhu-modularity input from
Theorem~\ref{v4-arith:thm:P2-resolution}; the conditional Petersson condition
supplies the positive Hermitian form needed for P3 in the restricted
arithmetic branch Corollary~\ref{v4-arith:kp:cor:P3-unconditional}.
H-P\(^{\mathrm{Ar}}\) supplies the self-adjoint operator, the
zero-divisor determinant, and the required Bost--Connes trace.
The specified (Reg) condition and operator comparison give
VB\(^{*}\) on the primary Fock; positivity, determinacy and the
required support are all retained.
Theorem~\ref{v4-arith:rh:thm:VB-implies-F-RH} gives the zero-placement clause
of F-RH. By Remark~\ref{v4-arith:rh:rem:equivalence-to-RH}, this clause is
RH.
\end{proof}

\begin{corollary}[A determinant--positivity sufficient condition for RH]
\label{v4-arith:platonic:cor:single-positivity}
Assume P1, P2 and P3 on their stated domains, together with the
specified H-P\(^{\mathrm{Ar}}\) determinant--trace realization.
Assume its arithmetic moment functional is positive and determinate,
its representing measure has the support required by (Reg), and
item~1 of
Hypothesis~\ref{v4-arith:hpvbconv:hyp:operator-quantitative-comparison}
identifies (Reg) with VB\(^{*}\) on this operator domain. Then RH holds.
RH by itself is not asserted to construct this realization,
support identification, or comparison.
\end{corollary}
\begin{proof}
The positive determinate moment functional and supplied support
give (Reg). The assumed operator comparison gives VB\(^{*}\).
The remaining P1--P3 and determinant--trace hypotheses now permit
Theorem~\ref{v4-arith:platonic:thm:RH-completion} to be applied.
\end{proof}
